\chapter{Roseola (HHV-6 / Sixth Disease)}
\index{Roseola (HHV-6 / Sixth Disease)}
\label{sec:Roseola (HHV-6 / Sixth Disease)}

\section{Overview}
Roseola infantum, also known as sixth disease or exanthem subitum, is a common viral exanthem of early childhood caused by human herpesvirus 6 (HHV-6). It is characterized by high fever lasting 3--5 days followed by a rose-colored maculopapular rash that appears after the fever defervesces. Nearly all children are infected by age 2--3 years, and most infections are uncomplicated.
\section{Classification}

\begin{description}[font=\normalfont\bfseries,leftmargin=3em,labelwidth=2.5em]
  \item[Cause] Infectious (Viral)
  \item[Organ System] Pediatric / Infectious
  \item[Pathophysiology] Primary infection with HHV-6 (or less commonly HHV-7) causing lytic infection of CD4+ T lymphocytes
  \item[Duration] Acute (5--7 days)
  \item[Onset] Sudden
  \item[Mode of Transmission] Respiratory droplets and saliva (virus persists lifelong)
  \item[Clinical Course] Acute, self-limited
  \item[Severity] Mild
  \item[Extent of Disease] Systemic
  \item[Age Group] Infants and Young Children (peak 6--15 months)
  \item[Epidemiology] Universal; >90\% of children seropositive by age 2
  \item[ICD-11 Code] 1D60.2
\end{description}

\section{Causes}
Primary infection with human herpesvirus 6 (HHV-6B) causes roseola. HHV-7 causes a similar but less common illness. The virus infects CD4+ T lymphocytes and persists in the body lifelong as a latent infection, like other herpesviruses. Transmission occurs through respiratory secretions and saliva from asymptomatic contacts. Reactivation can occur in immunocompromised individuals.

\section{Risk Factors}

\begin{itemize}[noitemsep]
  \item Age 6--15 months (peak incidence)
  \item Lack of prior exposure (all children are susceptible)
  \item Close contact with infected children
  \item Immunosuppression (risk of reactivation)
  \item Daycare attendance
  \item Siblings (increased exposure)
\end{itemize}

\section{Signs and Symptoms}

\subsection{Early symptoms}
\begin{itemize}[noitemsep]
  \item Early manifestations of roseola (hhv-6 / sixth disease) may be subtle and nonspecific
\end{itemize}

\subsection{Common symptoms}
\begin{itemize}[noitemsep]
  \item \subsection{Common symptoms}
  \item \textbf{Common symptoms:}
  \item High fever (39--40.5°C) lasting 3--5 days
  \item Fever resolves abruptly (crisis), followed by rash
  \item Rose-colored, blanching maculopapular rash beginning on the trunk and spreading to the face and extremities
  \item Rash typically lasts 1--3 days without desquamation
  \item Mild respiratory symptoms, cough, and cervical lymphadenopathy
  \item Nagayama spots (erythematous papules on the soft palate)
  \item Irritability and mild diarrhea
  \item Difficulty performing daily activities
\end{itemize}

\subsection{Less common symptoms}
\begin{itemize}[noitemsep]
  \item Atypical or less common presentations of roseola (hhv-6 / sixth disease) may occur
\end{itemize}

\subsection{Emergency warning signs}
\begin{itemize}[noitemsep]
  \item Severe or worsening symptoms of roseola (hhv-6 / sixth disease) require immediate medical evaluation
\end{itemize}
\section{Progression and Stages}
After an incubation period of 5--15 days, high fever develops abruptly without localizing signs. The child appears well despite high fever (fever-toxicity dissociation). Fever resolves abruptly on day 3--5, and within 12--24 hours the characteristic rash appears. The rash fades over 1--3 days. Full recovery is the rule without specific antiviral therapy.

\section{Complications}

\begin{itemize}[noitemsep]
  \item Febrile seizures (most common complication, in 5--15\% of cases)
  \item Meningitis and encephalitis (rare)
  \item Hepatitis and pneumonitis (rare)
  \item HHV-6 reactivation in immunocompromised patients (post-transplant)
  \item Reduced quality of life
  \item Emotional and psychological effects
\end{itemize}

\section{Diagnosis}

\begin{itemize}[noitemsep]
  \item Clinical diagnosis based on classic fever-rash pattern
  \item HHV-6 IgM serology or PCR for confirmation (rarely needed)
  \item Serology cannot distinguish primary infection from reactivation
  \item Complete blood count may show leukopenia and relative lymphocytosis
  \item Lumbar puncture if encephalitis suspected
\end{itemize}

\section{Prevention}

\begin{itemize}[noitemsep]
  \item Good hand hygiene
  \item No vaccine available
  \item Most children are infected asymptomatically, making prevention impractical
  \item Avoid exposure of immunocompromised children to known cases
  \item Go for regular check-ups so any problems are caught early
\end{itemize}

\section{Treatment Options}

\begin{itemize}[noitemsep]
  \item Supportive care: antipyretics (acetaminophen, ibuprofen) for fever
  \item Adequate hydration
  \item Anticonvulsant management for febrile seizures
  \item No specific antiviral therapy needed for immunocompetent children
  \item Ganciclovir or foscarnet for severe HHV-6 infections in immunocompromised patients
  \item Lifestyle changes and self-care
\end{itemize}

\section{Living with It}

\begin{itemize}[noitemsep]
  \item Follow your treatment plan as prescribed
  \item Keep up with regular appointments with your healthcare provider
  \item Adjust your daily habits as needed
  \item Reach out to family, friends, or support groups for help
  \item Keep learning about your condition and how to manage it
\end{itemize}

\section{Outlook}
The prognosis for roseola is excellent. It is a self-limited illness with full recovery in virtually all immunocompetent children. Febrile seizures are the most common complication and rarely have long-term consequences. Serious complications such as encephalitis are rare. HHV-6 remains latent lifelong but rarely causes disease in immunocompetent individuals after primary infection.
